\section{A Domain That Can Ask}
\label{sec:ch13-domain-that-can-ask}
\index{search subgraph!complete projection}%

Search is a small read model, not another owner for Session. Its project uses
the same runtime, federation package, paging support, and SQLite provider as
the existing services:

\begin{minted}[fontsize=\footnotesize]{xml}
<Project Sdk="Microsoft.NET.Sdk.Web">
  <PropertyGroup>
    <TargetFramework>net10.0</TargetFramework>
    <Nullable>enable</Nullable>
    <ImplicitUsings>enable</ImplicitUsings>
    <GenerateDocumentationFile>true</GenerateDocumentationFile>
    <NoWarn>$(NoWarn);CS1591</NoWarn>
  </PropertyGroup>
  <ItemGroup>
    <PackageReference Include="HotChocolate.ApolloFederation"
                      Version="16.6.1" />
    <PackageReference Include="HotChocolate.AspNetCore"
                      Version="16.6.1" />
    <PackageReference Include="HotChocolate.AspNetCore.CommandLine"
                      Version="16.6.1" />
    <PackageReference Include="HotChocolate.Data.EntityFramework"
                      Version="16.6.1" />
    <PackageReference Include="HotChocolate.Types.Analyzers"
                      Version="16.6.1" />
    <PackageReference Include="Microsoft.EntityFrameworkCore.Sqlite"
                      Version="10.0.11" />
  </ItemGroup>
</Project>
\end{minted}

The fourth local port is 5004:

\begin{minted}{json}
{
  "$schema": "https://json.schemastore.org/launchsettings.json",
  "profiles": {
    "http": {
      "commandName": "Project",
      "dotnetRunMessages": true,
      "launchBrowser": false,
      "applicationUrl": "http://localhost:5004",
      "environmentVariables": {
        "ASPNETCORE_ENVIRONMENT": "Development"
      }
    }
  }
}
\end{minted}

The read-model row carries everything needed to decide the page. The integer
key stays out of the schema; clients receive a Session entity instead.

\begin{minted}{csharp}
namespace Search;

public sealed record SessionSearchDocument(
    [property: GraphQLIgnore] int SessionId,
    DateTimeOffset StartsAt,
    double? AverageScore,
    int RatingCount);
\end{minted}

The seed is deliberately complete. Session 4 has no ratings, but it still has
a search document. A ratings table alone could not produce that row.

\begin{minted}{csharp}
namespace Search;

public static class SessionSearchData
{
    public static readonly
        IReadOnlyList<SessionSearchDocument> Documents =
    [
        new(1, At(9), 4.0, 3),
        new(2, At(10), 4.5, 2),
        new(3, At(11), 3.0, 1),
        new(4, At(14), null, 0)
    ];

    private static DateTimeOffset At(int hour) =>
        new(2026, 9, 14, hour, 0, 0, TimeSpan.Zero);
}
\end{minted}

The primary key is also the final ordering key. That turns equal start times or
equal scores into one deterministic order. It does not promise a snapshot
across separate requests; rows can still change between pages.

\begin{minted}{csharp}
using Microsoft.EntityFrameworkCore;

namespace Search;

public sealed class SearchContext(
    DbContextOptions<SearchContext> options) : DbContext(options)
{
    public DbSet<SessionSearchDocument> SessionSearch =>
        Set<SessionSearchDocument>();

    protected override void OnModelCreating(ModelBuilder model)
    {
        model.Entity<SessionSearchDocument>()
            .HasKey(s => s.SessionId);
        model.Entity<SessionSearchDocument>()
            .Property(s => s.StartsAt)
            .HasConversion(
                v => v.UtcDateTime,
                v => new DateTimeOffset(v, TimeSpan.Zero));
        model.Entity<SessionSearchDocument>()
            .HasData(SessionSearchData.Documents);
    }
}
\end{minted}

One enum is enough for the two orderings this example needs:

\begin{minted}{csharp}
namespace Search;

public enum SessionSearchOrder
{
    StartsAt,
    RatingDescending
}
\end{minted}

The resolver performs the important operations in one \code{IQueryable}.
Filtering and ordering happen before Hot Chocolate applies the paging
arguments. Both orders end with \code{SessionId}.

\begin{minted}{csharp}
using Microsoft.EntityFrameworkCore;

namespace Search;

[QueryType]
public static class Query
{
    [UsePaging(DefaultPageSize = 2, MaxPageSize = 10)]
    public static IQueryable<SessionSearchDocument> SearchSessions(
        SearchContext db,
        double? minimumScore,
        SessionSearchOrder? order)
    {
        var query = db.SessionSearch.AsNoTracking();
        if (minimumScore is not null)
        {
            query = query.Where(
                s => s.AverageScore >= minimumScore.Value);
        }
        return order is SessionSearchOrder.RatingDescending
            ? query.OrderByDescending(s => s.AverageScore)
                .ThenBy(s => s.SessionId)
            : query.OrderBy(s => s.StartsAt)
                .ThenBy(s => s.SessionId);
    }
}
\end{minted}

Notice what is absent: no call to Sessions, Ratings, or the router. Search can
decide the page locally because its projection already has the decision
values. That is the price and the point of the projection.
